% GENERATED FILE. Edit canonical lessons, then run npm run generate:books.
% canonical-source-hash: fnv1a64:2fb99f69bca31777
% canonical-lessons: PA-C248-diudhi, PA-C248-chaukhat, PA-C248-aral, PA-C248-sangal, PA-C248-khunti

\chapter{At the Door}
\label{ch:pa-a2-at-the-door}
\hlchaptermodality{248}

\begin{chapteropening}
\textbf{By the end of this chapter:} \emph{I can name the parts of a doorway.}

Complete the last lesson of chapter 248: I can name the parts of a doorway.
\end{chapteropening}

\section[\texorpdfstring{\pa{ਡਿਉਢੀ} (ḍiuḍhī)}{ਡਿਉਢੀ (ḍiuḍhī)}]{\pa{ਡਿਉਢੀ} (ḍiuḍhī) — a covered entrance; a threshold}

\label{lesson:PA-C248-diudhi}



\begin{quote}
Before the new one: say the Punjabi for a low stool, then the Punjabi for a trunk.
\end{quote}

\subsection*{You'll want to know: \pa{ਡਿਉਢੀ}}
\textbf{\pa{ਡਿਉਢੀ}} — \emph{ḍiuḍhī} — \textquotedblleft{}a covered entrance; a threshold\textquotedblright{}.

\begin{quote}
\pa{ਅਸੀਂ} \pa{ਜੁੱਤੀਆਂ} \pa{ਡਿਉਢੀ} \pa{ਵਿੱਚ} \pa{ਰੱਖਦੇ} \pa{ਹਾਂ।}
\end{quote}

\begin{quote}
\emph{asīṁ juttīāṁ ḍiuḍhī vich rakkhde hāṁ.}
\end{quote}

\begin{quote}
We keep our shoes in the covered entrance.
\end{quote}

A feminine noun.

\begin{cousinweb}
\textbf{\pa{ਡਿਉਢੀ}} (\emph{ḍiuḍhī}) is inherited from a reconstructed \emph{*dēhuḍī}, \textquotedblleft{}mound, threshold\textquotedblright{}; CDIAL lists Punjabi \emph{ḍeuḍhī}, \textquotedblleft{}threshold\textquotedblright{} (CDIAL 6565).
\end{cousinweb}

\begin{grammarlens}[title={What you've built}]
One more word for this chapter.
\end{grammarlens}

\subsection*{Your turn}
\begin{itemize}
\raggedright
  \item \emph{Say it:} \emph{ḍiuḍhī}
  \item \emph{Say it:} \emph{ḍiuḍhī}, once more
  \item \emph{Say it:} \emph{asīṁ juttīāṁ ḍiuḍhī vich rakkhde hāṁ.} — We keep our shoes in the covered entrance.
  \item \emph{Recall:} say the Punjabi for a low stool, then the Punjabi for a trunk, then say \emph{ḍiuḍhī} again
\end{itemize}

\begin{tcolorbox}[breakable,colback=teal!4,colframe=teal!35!black,title={Before you move on}]
What does \pa{ਡਿਉਢੀ} mean? (\textquotedblleft{}A covered entrance; a threshold\textquotedblright{}.) Say it once more.
\end{tcolorbox}
\section[\texorpdfstring{\pa{ਚੌਖਟ} (chaukhaṭ)}{ਚੌਖਟ (chaukhaṭ)}]{\pa{ਚੌਖਟ} (chaukhaṭ) — a door frame}

\label{lesson:PA-C248-chaukhat}



\begin{quote}
Before the new one: say the Punjabi for a trunk, then the Punjabi for a covered entrance.
\end{quote}

\subsection*{You'll want to know: \pa{ਚੌਖਟ}}
\textbf{\pa{ਚੌਖਟ}} — \emph{chaukhaṭ} — \textquotedblleft{}a door frame\textquotedblright{}.

\begin{quote}
\pa{ਚੌਖਟ} \pa{ਪੁਰਾਣੀ} \pa{ਹੈ} \pa{ਪਰ} \pa{ਦਰਵਾਜ਼ਾ} \pa{ਨਵਾਂ} \pa{ਹੈ।}
\end{quote}

\begin{quote}
\emph{chaukhaṭ purāṇī hai par darvāzā navāṁ hai.}
\end{quote}

\begin{quote}
The door frame is old, but the door is new.
\end{quote}

A feminine noun.

\begin{cousinweb}
\textbf{\pa{ਚੌਖਟ}} (\emph{chaukhaṭ}) is inherited from a reconstructed \emph{*catuṣkāṣṭha}, \textquotedblleft{}rectangular frame\textquotedblright{}; CDIAL lists Punjabi \emph{caukhaṭ}, \textquotedblleft{}doorframe\textquotedblright{} (CDIAL 4631).
\end{cousinweb}

\begin{grammarlens}[title={What you've built}]
One more word for this chapter.
\end{grammarlens}

\subsection*{Your turn}
\begin{itemize}
\raggedright
  \item \emph{Say it:} \emph{chaukhaṭ}
  \item \emph{Say it:} \emph{chaukhaṭ}, once more
  \item \emph{Say it:} \emph{chaukhaṭ purāṇī hai par darvāzā navāṁ hai.} — The door frame is old, but the door is new.
  \item \emph{Recall:} say the Punjabi for a trunk, then the Punjabi for a covered entrance, then say \emph{chaukhaṭ} again
  \item \emph{Write it:} copy \textbf{\pa{ਚੌਖਟ}} once, checking the whole word
\end{itemize}

\begin{tcolorbox}[breakable,colback=teal!4,colframe=teal!35!black,title={Before you move on}]
What does \pa{ਚੌਖਟ} mean? (\textquotedblleft{}A door frame\textquotedblright{}.) Say it once more.
\end{tcolorbox}
\section[\texorpdfstring{\pa{ਅਰਲ} (aral)}{ਅਰਲ (aral)}]{\pa{ਅਰਲ} (aral) — a door bolt}

\label{lesson:PA-C248-aral}



\begin{quote}
Before the new one: say the Punjabi for a covered entrance, then the Punjabi for a door frame.
\end{quote}

\subsection*{You'll want to know: \pa{ਅਰਲ}}
\textbf{\pa{ਅਰਲ}} — \emph{aral} — \textquotedblleft{}a door bolt\textquotedblright{}.

\begin{quote}
\pa{ਰਾਤ} \pa{ਨੂੰ} \pa{ਦਰਵਾਜ਼ੇ} \pa{ਦੀ} \pa{ਅਰਲ} \pa{ਬੰਦ} \pa{ਹੁੰਦੀ} \pa{ਹੈ।}
\end{quote}

\begin{quote}
\emph{rāt nū̃ darvāze dī aral band hundī hai.}
\end{quote}

\begin{quote}
At night the door bolt is shut.
\end{quote}

A feminine noun.

\begin{cousinweb}
\textbf{\pa{ਅਰਲ}} (\emph{aral}) is inherited from Sanskrit \emph{ararī}, \textquotedblleft{}leaf of a door\textquotedblright{}; CDIAL lists Punjabi \emph{aral}, \textquotedblleft{}wooden bolt of a door\textquotedblright{} (CDIAL 605).
\end{cousinweb}

\begin{grammarlens}[title={What you've built}]
One more word for this chapter.
\end{grammarlens}

\subsection*{Your turn}
\begin{itemize}
\raggedright
  \item \emph{Say it:} \emph{aral}
  \item \emph{Say it:} \emph{aral}, once more
  \item \emph{Say it:} \emph{rāt nū̃ darvāze dī aral band hundī hai.} — At night the door bolt is shut.
  \item \emph{Recall:} say the Punjabi for a covered entrance, then the Punjabi for a door frame, then say \emph{aral} again
  \item \emph{Write it:} copy \textbf{\pa{ਅਰਲ}} once, checking the whole word
\end{itemize}

\begin{tcolorbox}[breakable,colback=teal!4,colframe=teal!35!black,title={Before you move on}]
What does \pa{ਅਰਲ} mean? (\textquotedblleft{}A door bolt\textquotedblright{}.) Say it once more.
\end{tcolorbox}
\section[\texorpdfstring{\pa{ਸੰਗਲ} (saṅgal)}{ਸੰਗਲ (saṅgal)}]{\pa{ਸੰਗਲ} (saṅgal) — a chain}

\label{lesson:PA-C248-sangal}



\begin{quote}
Before the new one: say the Punjabi for a door frame, then the Punjabi for a door bolt.
\end{quote}

\subsection*{You'll want to know: \pa{ਸੰਗਲ}}
\textbf{\pa{ਸੰਗਲ}} — \emph{saṅgal} — \textquotedblleft{}a chain\textquotedblright{}.

\begin{quote}
\pa{ਦਰਵਾਜ਼ੇ} \pa{ਉੱਤੇ} \pa{ਇੱਕ} \pa{ਭਾਰੀ} \pa{ਸੰਗਲ} \pa{ਹੈ।}
\end{quote}

\begin{quote}
\emph{darvāze utte ikk bhārī saṅgal hai.}
\end{quote}

\begin{quote}
There is a heavy chain on the door.
\end{quote}

A masculine noun.

\begin{cousinweb}
\textbf{\pa{ਸੰਗਲ}} (\emph{saṅgal}) is inherited from Sanskrit \emph{śṛṅkhala}, \textquotedblleft{}chain\textquotedblright{}; CDIAL lists Punjabi \emph{saṅgal}, \textquotedblleft{}chain\textquotedblright{} (CDIAL 12580).
\end{cousinweb}

\begin{grammarlens}[title={What you've built}]
One more word for this chapter.
\end{grammarlens}

\subsection*{Your turn}
\begin{itemize}
\raggedright
  \item \emph{Say it:} \emph{saṅgal}
  \item \emph{Say it:} \emph{saṅgal}, once more
  \item \emph{Say it:} \emph{darvāze utte ikk bhārī saṅgal hai.} — There is a heavy chain on the door.
  \item \emph{Recall:} say the Punjabi for a door frame, then the Punjabi for a door bolt, then say \emph{saṅgal} again
  \item \emph{Write it:} copy \textbf{\pa{ਸੰਗਲ}} once, checking the whole word
\end{itemize}

\begin{tcolorbox}[breakable,colback=teal!4,colframe=teal!35!black,title={Before you move on}]
What does \pa{ਸੰਗਲ} mean? (\textquotedblleft{}A chain\textquotedblright{}.) Say it once more.
\end{tcolorbox}
\section[\texorpdfstring{\pa{ਖੂੰਟੀ} (khūṇṭī)}{ਖੂੰਟੀ (khūṇṭī)}]{\pa{ਖੂੰਟੀ} (khūṇṭī) — a wall peg}

\label{lesson:PA-C248-khunti}



\begin{quote}
Before the new one: say the Punjabi for a door bolt, then the Punjabi for a chain.
\end{quote}

\subsection*{You'll want to know: \pa{ਖੂੰਟੀ}}
\textbf{\pa{ਖੂੰਟੀ}} — \emph{khūṇṭī} — \textquotedblleft{}a wall peg\textquotedblright{}.

\begin{quote}
\pa{ਮੇਰਾ} \pa{ਬਸਤਾ} \pa{ਖੂੰਟੀ} \pa{ਉੱਤੇ} \pa{ਹੈ।}
\end{quote}

\begin{quote}
\emph{merā bastā khūṇṭī utte hai.}
\end{quote}

\begin{quote}
My school bag is on the peg.
\end{quote}

A feminine noun.

\begin{cousinweb}
\textbf{\pa{ਖੂੰਟੀ}} (\emph{khūṇṭī}) goes back to a reconstructed \emph{*khuṭṭa}, \textquotedblleft{}peg, post\textquotedblright{}, as a word taken from Hindi: CDIAL lists Punjabi \emph{khū̃ṭī}, \textquotedblleft{}stake, peg\textquotedblright{}, as borrowed from Hindi \emph{khū̃ṭī}, \textquotedblleft{}small peg\textquotedblright{} (CDIAL 3893).
\end{cousinweb}

\begin{grammarlens}[title={What you've built}]
One more word for this chapter.
\end{grammarlens}

\subsection*{Your turn}
\begin{itemize}
\raggedright
  \item \emph{Say it:} \emph{khūṇṭī}
  \item \emph{Say it:} \emph{khūṇṭī}, once more
  \item \emph{Say it:} \emph{merā bastā khūṇṭī utte hai.} — My school bag is on the peg.
  \item \emph{Recall:} say the Punjabi for a door bolt, then the Punjabi for a chain, then say \emph{khūṇṭī} again
  \item \emph{Write it:} copy \textbf{\pa{ਖੂੰਟੀ}} once, checking the whole word
\end{itemize}

\begin{tcolorbox}[breakable,colback=teal!4,colframe=teal!35!black,title={Before you move on}]
What does \pa{ਖੂੰਟੀ} mean? (\textquotedblleft{}A wall peg\textquotedblright{}.) Say it once more.
\end{tcolorbox}
